\begin{titlepage}
\centering
{\Large Notas didácticas de revisión técnica\par}
\vspace{1.0cm}
{\huge\bfseries Kimi Linear, Sparse Attention y la arquitectura de modelos de la próxima generación}\par
\vspace{0.5cm}
{\Large Zhang Xiaojun conversa con Yang Songlin (杨松琳)\par}
\vspace{0.35cm}
{\large Elaborado a partir del video público de Zhang Xiaojun Podcast, la transcripción generada en GPU, la estructura de capítulos y diagramas conceptuales propios\par}
\vspace{0.3cm}
{\large 2025-11-03\par}
\vspace{0.85cm}
\includegraphics[width=0.88\textwidth,height=0.42\textheight,keepaspectratio]{cover.jpg}\par
\vfill
\begin{tcolorbox}[width=0.92\textwidth, colback=black!2!white, colframe=black!55, sharp corners]
\textbf{Título del video}: 119. Kimi Linear, ¿Minimax M2? Yang Songlin repasa la historia de las variantes algorítmicas y anticipa propuestas de mejora para las arquitecturas futuras\par
\textbf{Canal del video}: Zhang Xiaojun Podcast\par
\textbf{Duración del video}: 01:43:26\par
\textbf{Enlace del video}: \href{\videourl}{\nolinkurl{\videourl}}\par
\textbf{Notas sobre el material}: YouTube no ofrece subtítulos disponibles; el texto se elaboró a partir de la transcripción con faster-whisper large-v3 en CUDA, los capítulos del video, la portada y figuras didácticas propias. Las tomas fijas de la entrevista no se reutilizan en el cuerpo del texto.\par
\textbf{Página del proyecto}: \href{\repourl}{\nolinkurl{\repourl}}\par
\end{tcolorbox}
\end{titlepage}

\tableofcontents
\newpage
